\documentclass[11pt]{article}

\usepackage[margin=25mm]{geometry}

\usepackage{amsmath,amssymb,fontspec,unicode-math}

\IfFontExistsTF{DejaVu Serif}{\setmainfont{DejaVu Serif}}{\setmainfont{Latin Modern Roman}}

\IfFontExistsTF{Latin Modern Math}{\setmathfont{Latin Modern Math}}{\setmathfont{DejaVu Math TeX Gyre}}

\usepackage{xurl,hyperref}

\setlength{\parindent}{0pt}\setlength{\parskip}{6pt}

\begin{document}

\title{RDD: Retrieval-Based Demonstration Decomposer for Planner Alignment in Long-Horizon Tasks}

\author{ELVIA preliminary AI analysis}\date{21 October 2026}\maketitle

Preliminary assessment: Major.\par

Independent scientific review remains pending. Author adoption has not been requested.

The current paper and author release still omit the shortest feasible dynamic-programming prefix, contradicting the stated complete-partition optimum guarantee on a legal four-frame example. A separate Minor concern distinguishes linear scoring calls from the implementation’s quadratic length checks and path copies. Robot success rates are not independently refuted.

\section*{Review scope}

Latest-version, claim-focused review of the cited author-paper passages and applicable current author-release paths. The stated counterexamples and scope controls were checked analytically; no author program, old verifier, training run or benchmark was executed.

\section*{Problem 1: The first feasible prefix is skipped}

Source: Corollary 3.1.1, PDF p. 4; Algorithm 1 and Appendix A.2, pp. 22–23; current rdd/algorithms.py, max\_sum\_partition and rdd\_score

Assessment: supported; impact: major

\subsection*{Source claim (paraphrase)}

Under additive interval scores and allowed lengths, Algorithm 1 is said to find an optimal full partition.

\subsection*{Analysis}

Take \(N=4\), minimum length \(2\) and maximum length \(4\). All frame features can equal \((1,0)\); a single two-frame prior interval has duration \(1\), so its endpoint feature matches every candidate. With the author scorer’s default mode, \(\alpha=1\) and \(\beta=0\), angular distance is zero and its score is the negative frame count times \(|1-(e-s)|\). Two-frame segments score \(0\), while the four-frame segment scores \(-8\). The loop begins at \(3\), leaving \(D_2=-\infty\), and hence returns \(-8\) rather than the feasible two-segment optimum \(0\). Interpreting the paper’s outside duration as \(e-s\) instead gives the full segment \(-6\), with the same strict contradiction. If the maximum length is \(3\), the valid two-segment solution is still missed. Invalid-edge scores of \(-\infty\) do not restore the omitted feasible state.

\subsection*{Finding and scope}

The listed algorithm and matching author release fail the stated universal optimum guarantee. Standard additive dynamic programming remains valid; no robot success-rate effect is inferred.

\subsection*{Proposed correction}

Start at the shortest feasible prefix, preserve genuinely infeasible states and backtrack a full partition. The corrected standard recurrence is distinct from a new robot experiment.

The latest target arXiv version and freshly retrieved proceedings Algorithm 1 retain the omitted prefix; the current author repository retains the same loop lower bound. The existing standard-recurrence correction is analytic, with no new robot execution or author adoption established.

\section*{Problem 2: Linear scorer calls do not bound total implementation cost}

Source: Corollary 3.1.1, PDF p. 4; speed discussion, PDF pp. 8–9; Appendix A.2, PDF p. 23.

Assessment: supported; impact: minor

\subsection*{Source claim (paraphrase)}

With bounded segment length, the paper counts linear scoring calls and later describes linear algorithm runtime.

\subsection*{Analysis}

With allowed lengths \(2\) through \(4\), the release checks all previous breakpoints before discarding illegal lengths. It makes \(N(N+1)/2+N-5=\Theta(N^2)\) length checks but only \(3N-7\) scorer calls. Full-prefix partition copies add further work. The finite Table 14 budget \(44549/386\simeq115.41\,\mathrm s\) is a valid call-count calculation, not a total asymptotic-runtime proof.

\subsection*{Finding and scope}

The call-count result is preserved; an asymptotic linear total-runtime claim needs restricted candidate loops and path storage assumptions. Finite empirical timings are not refuted.

\subsection*{Proposed correction}

Restrict candidate breakpoints to the valid length band and store predecessor links. With declared interval-access and scoring cost \(C_J\), that gives \(O(NW(1+C_J))\) time for width \(W\).

The current author function still scans every previous breakpoint before checking lengths and copies full partition prefixes. The bounded scoring-call statement remains valid; no new end-to-end timing was measured.

\section*{Limits}

Independent scientific adjudication remains pending.\par

No author code, model training or full benchmark evaluation was executed in this revision.\par

Version-specific source defects do not establish which code produced the published experiments.\par

Corrected benchmark scores, empirical effect sizes and author adoption remain unverified unless explicitly stated.\par

Local contradictions and specification differences do not establish the magnitude or direction of changes to the published benchmark scores.\par

No author contact or adoption; novelty and complete historical-run provenance are unestablished.\par

\section*{Source identity}

NeurIPS 2025 official proceedings PDF (current publisher bytes, checked 2026-10-21)

\url{https://papers.nips.cc/paper_files/paper/2025/file/43ae0b566b802b2d00e525b672794b82-Paper-Conference.pdf}

Reviewed PDF SHA-256: \nolinkurl{b6609ce921186e9eeaffe43d2d0f5cbc3d6c1f6213fe8812ea9371f5c103a0f6}

This is the same research work as the predecessor, matched by title, author identity and official publisher channel. Its freshly retrieved exact PDF bytes and real retrieval time are recorded. Current target arXiv metadata and the relevant available author-version passages were checked separately; background works were excluded. Original report/source history is unchanged.

\end{document}
